% id: 211-443
% book: 개념원리 기하 · 25 벡터의 내적과 두 벡터가 이루는 각
% section: 연습문제 STEP 2
% figure: none
% answer: $\dfrac{2}{3}$
% answer_source: 답지
% variant_level: 0 (원본 전사)
% 이 파일은 build.mjs 가 items.json 에서 만든다. 고칠 때는 items.json 을 고친다.
\smash{\raisebox{1.5mm}{\dmkichul{개념원리 기하 211쪽 443번}}}
좌표공간의 세 점~$\pt{O}(0{,}\mkern3mu\mathopen{}0{,}\mkern3mu\mathopen{}0)$, $\pt{A}(2{,}\mkern3mu\mathopen{}3{,}\mkern3mu\mathopen{}-1)$, $\pt{B}(-1{,}\mkern3mu\mathopen{}2{,}\mkern3mu\mathopen{}2)$에 대하여 벡터 $\overrightarrow{\pt{OA}}$의 직선~$\pt{OB}$ 위로의 정사영의 크기를 구하시오.
